\section{Representación visual general: de la inducción humana con pocos ejemplos a un objetivo medible}

Antes de pasar al modelo, hay que definir el objetivo. Una representación que solo funciona dentro del espacio de etiquetas del preentrenamiento se parece más a un conjunto de características especializadas. Una \term{general visual representation} (representación visual general) debería permitir que una tarea nueva alcance un desempeño confiable con muy pocos ejemplos y un costo de adaptación reducido. También debería abarcar problemas visuales distintos: imágenes naturales, escenas estructuradas, percepción remota, conteo y geometría.

\subsection{El objetivo no es «ver», sino adaptarse rápidamente}

El ponente plantea primero una visión centrada en los usuarios reales: en el futuro, una persona común podrá enseñarle al sistema una tarea visual nueva con unos cuantos ejemplos, sin necesidad de programar. Lo esencial no es la apariencia de robot, sino la adaptation interface: cuando el usuario proporciona ejemplos, ¿puede el sistema aprovechar la representación existente para entender rápidamente una nueva regla de clasificación?

\lecturefigure{slide-02-goal-general-representation.jpg}{El objetivo de la investigación es una general visual representation, no la puntuación máxima en un solo benchmark.}{Video oficial 00:36.}

\begin{knowledgebox}{Lectura de la figura: la representation es un activo común previo a la tarea}
El modelo preentrenado puede escribirse como un mapeo de características $h=f_{\theta}(x)$, y la tarea posterior aprende después $g_{\phi}(h)$. Si $f_{\theta}$ es realmente general, una nueva $g_{\phi}$ debería requerir solo pocos datos o una actualización pequeña de parámetros. La figura no especifica clasificación, detección ni segmentación precisamente para subrayar que la representación antecede a cualquier tarea concreta. Para evaluarla hay que observar un conjunto de tareas posteriores, no solo la upstream accuracy.
\end{knowledgebox}

\lecturefigure{slide-03-why-visual-representation.jpg}{Una buena representación visual debería dar un «kick-start» a muchas tareas con entrada visual y reducir el costo de enseñanza para quienes no son especialistas.}{Video oficial 01:40.}

\begin{importantbox}{Qué significa kick-start}
Kick-start no significa resolver la tarea a costo cero, sino reducir de manera notable tres tipos de costo: \textbf{sample cost} (cuántos ejemplos etiquetados se necesitan), \textbf{optimization cost} (cuánto tiempo hay que entrenar y cuántos parámetros hay que actualizar) y \textbf{specification cost} (con cuánta precisión debe describir la tarea el usuario). Una representación puede mejorar la precisión final sin reducir estos costos y, por lo tanto, no puede considerarse fácil de adaptar.
\end{importantbox}

\teachervoice{Lucas explica la visión de la investigación con la idea de que «un niño o sus padres puedan enseñarle a un robot sencillo», y aclara que la representación visual es solo una parte de un sistema inteligente completo. Nota de clase: expresar el valor del aprendizaje de representaciones como costo de adaptación para el usuario es más verificable que afirmar de forma general que «el modelo es más inteligente».}

\subsection{Qué referencia ofrece la representación visual humana}

Las personas no solo son buenas para reconocer objetos conocidos. Con flores, imágenes satelitales y objetos abstractos, el ponente muestra que, aunque la semántica de las categorías sea desconocida, una persona puede inferir la regla a partir de unos pocos ejemplos. Lo que vale la pena imitar es la capacidad de combinar elementos de la representación para formar nuevas fronteras de decisión, no la de memorizar una enorme cantidad de nombres de categorías. Al leer los tres grupos de figuras siguientes, hay que determinar en cada caso si la tarea depende de experiencia semántica de largo plazo, de estadísticas visuales de bajo nivel o de reglas abstractas como la cantidad y las relaciones. Solo una representación que se adapte rápidamente a través de todos estos mecanismos se acerca a ser «general».

\lecturefigure{slide-04-human-visual-representation.jpg}{La representación visual humana se usa como referencia de la capacidad de inducción con pocos ejemplos.}{Video oficial 01:48.}

\begin{knowledgebox}{Lectura de la figura: el éxito con pocos ejemplos puede provenir de mecanismos distintos}
La tarea de flores puede aprovechar la experiencia visual de largo plazo; las imágenes satelitales requieren transferir texturas y estructura espacial; la tarea de objetos abstractos puede depender de reglas de conteo o de relaciones. Las tres se presentan como few-shot classification, pero requieren invariance distintas. Un benchmark que solo contenga objetos naturales similares no puede distinguir entre «memorizar plantillas semánticas» y «formar una estructura visual general».
\end{knowledgebox}

\lecturefigure{slide-05-vtab-flowers.jpg}{Las categorías de flores del ejemplo de VTAB están cerca de la distribución de imágenes naturales; con pocos ejemplos basta para formar el concepto de cada categoría.}{Video oficial 02:00; Zhai et al. 2019.}

\begin{knowledgebox}{Lectura de la figura: primero hay que juzgar la distancia entre la tarea y la distribución de preentrenamiento}
Las imágenes de flores tienen texturas, colores y contornos reales, y suelen estar cerca de un preentrenamiento del tipo ImageNet. Si el modelo rinde bien en este tipo de tarea, puede deberse a una representación transferible o a la superposición semántica con las categorías o imágenes upstream. Al leer un benchmark, el primer paso no es mirar la puntuación, sino preguntar qué tan cerca están downstream y upstream en cuanto a origen de las imágenes, semántica de las etiquetas y estadísticas visuales.
\end{knowledgebox}

\lecturefigure{slide-06-vtab-objects.jpg}{La tarea de objetos abstractos exige inducir la regla de cada categoría a partir de la cantidad, la forma o las relaciones.}{Video oficial 03:04; VTAB structured tasks.}

\begin{warningbox}{Few-shot no es tan simple como «entrenar con solo cinco imágenes»}
Que el modelo vea pocos ejemplos etiquetados en la tarea posterior no significa que no haya visto patrones visuales similares. Hay que distinguir pretraining exposure, downstream labeled examples, hyperparameter tuning data y test examples. Si la recipe se ajusta repetidamente sobre varias tareas posteriores, el sistema en conjunto usa mucha más información que los $k$ ejemplos de una sola tarea.
\end{warningbox}

\subsection{Resumen de la sección}

Lo central de una representación general es adaptarse rápidamente a reglas diversas, no la precisión en una sola clasificación. Los ejemplos humanos indican que hay que cubrir a la vez semántica conocida, domain shift y razonamiento estructurado; esto conduce directamente al marco de medición multitarea de VTAB.
